% =============================================================================
% Le projet AG expliqué depuis zéro (EVCSP, groupe 1)
% Compilation : lualatex explication_du_sujet.tex (deux fois, pour la table
% des matières)
% =============================================================================
\documentclass[11pt,a4paper]{article}

\usepackage{fontspec}
\usepackage[french]{babel}
\usepackage[a4paper,margin=2.3cm]{geometry}
\usepackage{amsmath,amssymb}
\usepackage{booktabs,tabularx,array}
\usepackage{xcolor}
\usepackage{enumitem}
\usepackage{tikz}
\usetikzlibrary{arrows.meta,positioning,shapes.geometric,calc}
\usepackage[most]{tcolorbox}
\usepackage[colorlinks=true,linkcolor=blue!50!black,urlcolor=blue!50!black]{hyperref}

\setlength{\parindent}{0pt}
\setlength{\parskip}{0.55em}
\setlist{itemsep=0.15em,topsep=0.3em}
\renewcommand{\arraystretch}{1.25}

% Encadrés
\newtcolorbox{idee}[1][]{colback=blue!4,colframe=blue!45!black,fonttitle=\bfseries,title={#1},sharp corners,boxrule=0.6pt}
\newtcolorbox{attention}[1][]{colback=orange!6,colframe=orange!70!black,fonttitle=\bfseries,title={#1},sharp corners,boxrule=0.6pt}

% Raccourcis
\newcommand{\As}{\mathbf{A}^{*}}
\newcommand{\Bs}{\mathbf{B}^{*}}
\newcommand{\Ds}{\mathbf{D}^{*}}
\newcommand{\Cm}{\mathbf{C}}
\newcommand{\Q}{\mathbf{Q}}
\newcommand{\Sm}{\mathbf{S}}
\newcommand{\code}[1]{\texttt{#1}}

\title{\textbf{Le projet AG expliqué depuis zéro}\\[0.3em]
\large Optimisation d'un stratifié carbone/époxyde par algorithme génétique\\
et méthode polaire --- EVCSP, groupe 1}
\author{}
\date{}

\begin{document}
\maketitle
\vspace{-3em}

\begin{idee}[Comment lire ce document]
Il explique \textbf{tout le sujet} comme si tu ne l'avais jamais vu, dans
l'ordre logique : d'abord \emph{ce qu'on fabrique} (sections 1 à 3),
ensuite \emph{comment on décrit sa rigidité} (section 4), puis \emph{ce qu'on
cherche} (section 5) et enfin \emph{comment on le cherche} (section 6). Les
sections 7 et 8 font le lien avec le code et les résultats. Les numéros de
slides renvoient au cours (P1, P2) et à l'énoncé.
\end{idee}

\tableofcontents
\newpage

% =============================================================================
\section{Le projet en une phrase}
% =============================================================================
On cherche dans quelles directions orienter les \textbf{20 couches de fibres
de carbone} d'une plaque pour qu'elle se comporte mécaniquement d'une façon
précise (imposée par l'énoncé), et on fait chercher ces 20 angles par un
\textbf{algorithme génétique}, un programme qui imite l'évolution naturelle.

% =============================================================================
\section{Qu'est-ce qu'un composite stratifié ?}
% =============================================================================
\subsection{Le pli}
Un \textbf{pli} (ou couche) est une feuille très fine, ici
$t_{\text{ply}} = 0{,}125$~mm, faite de :
\begin{itemize}
  \item \textbf{fibres de carbone} toutes parallèles : très rigides et très
        résistantes \emph{dans leur direction} ;
  \item \textbf{résine époxyde} (la \og matrice\fg) qui colle les fibres :
        molle.
\end{itemize}
Conséquence : un pli est \textbf{très rigide dans le sens des fibres} et
\textbf{mou en travers}.

\begin{center}
\begin{tabularx}{\linewidth}{l l X}
\toprule
Grandeur & Valeur & Ce que ça veut dire\\
\midrule
$E_1$ & 161\,000 MPa & rigidité en tirant \textbf{dans le sens des fibres}
  (l'acier est à 210\,000 MPa)\\
$E_2$ & 9\,000 MPa & rigidité en tirant \textbf{en travers} : 18 fois moins !\\
$G_{12}$ & 6\,100 MPa & rigidité en \textbf{cisaillement} (on fait
  \og glisser\fg{} le pli sur lui-même)\\
$\nu_{12}$ & 0,26 & coefficient de Poisson : si on allonge le pli de
  $\varepsilon$ selon 1, il rétrécit de $0{,}26\,\varepsilon$ selon 2\\
\bottomrule
\end{tabularx}
\end{center}

Un matériau dont les propriétés dépendent de la direction est dit
\textbf{anisotrope}. Le pli est un cas particulier : il a deux directions
privilégiées perpendiculaires (fibres et travers), il est
\textbf{orthotrope}.

\subsection{L'orientation d'un pli : les directions 1, 2 et \texorpdfstring{$x$, $y$, $z$}{x, y, z}}
On utilise \textbf{deux repères} différents, et il ne faut pas les
confondre :
\begin{itemize}
  \item le repère \textbf{de la plaque} $(x, y, z)$ : fixe, le même pour tous
        les plis. $x$ et $y$ sont dans le plan de la plaque, $z$ est dans
        l'épaisseur (vers le haut) ;
  \item le repère \textbf{du pli} $(1, 2)$ : propre à chaque pli, il tourne
        avec lui. \textbf{1 = direction des fibres}, \textbf{2 = direction
        perpendiculaire aux fibres}, dans le plan.
\end{itemize}
L'angle $\delta$ d'un pli est l'angle entre l'axe $x$ de la plaque et ses
fibres (l'axe 1), compté dans le sens trigonométrique :

\begin{center}
\begin{tikzpicture}[>={Stealth[length=2.6mm]},scale=1.0]
  % --- Left: plate frame and one ply rotated by delta ---
  \begin{scope}
    % ply (square) with fibres rotated by 30 deg
    \fill[blue!5] (-2,-2) rectangle (2,2);
    \draw[gray] (-2,-2) rectangle (2,2);
    \begin{scope}
      \clip (-2,-2) rectangle (2,2);
      \foreach \k in {-6,...,6}
        \draw[blue!35,thin,rotate=30] (-4,\k*0.35) -- (4,\k*0.35);
    \end{scope}
    % plate axes x, y
    \draw[->,thick] (0,0) -- (2.7,0) node[right] {$x$ (plaque)};
    \draw[->,thick] (0,0) -- (0,2.7) node[above] {$y$ (plaque)};
    % ply axes 1, 2
    \draw[->,very thick,red!75!black] (0,0) -- (30:2.4) node[above right] {1 (fibres)};
    \draw[->,very thick,red!75!black] (0,0) -- (120:2.2) node[above left] {2 ($\perp$ fibres)};
    % angle delta
    \draw[->] (0.9,0) arc (0:30:0.9);
    \node at (15:1.25) {$\delta$};
    \node[below,align=center,font=\small] at (0,-2.15) {Vue de dessus d'un pli : repère de la\\plaque $(x,y)$ et repère du pli $(1,2)$};
  \end{scope}
  % --- Right: examples of delta ---
  \begin{scope}[xshift=7.2cm,yshift=0.6cm]
    \foreach \a/\y in {0/1.4,90/0,45/-1.4}{
      \fill[blue!5] (-0.7,\y-0.5) rectangle (0.7,\y+0.5);
      \draw[gray] (-0.7,\y-0.5) rectangle (0.7,\y+0.5);
      \begin{scope}
        \clip (-0.7,\y-0.5) rectangle (0.7,\y+0.5);
        \foreach \k in {-4,...,4}
          \draw[blue!45,shift={(0,\y)},rotate=\a] (-1.5,\k*0.22) -- (1.5,\k*0.22);
      \end{scope}
      \node[right,font=\small] at (0.9,\y) {$\delta = \a^\circ$};
    }
    \node[right,font=\small] at (2.2,1.4) {fibres selon $x$};
    \node[right,font=\small] at (2.2,0) {fibres selon $y$};
    \node[right,font=\small] at (2.2,-1.4) {fibres en diagonale};
    \draw[->] (-1.1,-2.25) -- (-0.4,-2.25) node[right,font=\small] {$x$};
  \end{scope}
\end{tikzpicture}
\end{center}

Les constantes du pli ($E_1$, $E_2$, $G_{12}$, $\nu_{12}$) sont données
\textbf{dans le repère du pli} $(1,2)$. Les rigidités de la plaque
($\mathbf{A}$, $\mathbf{B}$, $\mathbf{D}$) sont écrites \textbf{dans le
repère de la plaque} $(x,y)$. Passer de l'un à l'autre, c'est
\og tourner le pli de $\delta$\fg{} (section 3.3).

Pourquoi l'intervalle $]-90^\circ, 90^\circ]$ ? Une fibre n'a pas de sens :
une fibre à $100^\circ$ est la même qu'une fibre à $-80^\circ$. Au pas de
$1^\circ$, il y a donc exactement \textbf{180 orientations} :
$-89, -88, \dots, 0, \dots, 89, 90$.

\subsection{Le stratifié}
Un \textbf{stratifié} est un empilement de plis collés, noté du bas vers le
haut : $[0/45/{-45}/90]$ = 4 plis, le premier à $0^\circ$, le dernier à
$90^\circ$. Chez nous : $n = 20$ plis, épaisseur totale
$h = n\,t_{\text{ply}} = 2{,}5$~mm.

\begin{idee}[L'idée clé du projet]
En choisissant les angles, on \og fabrique\fg{} un matériau sur mesure. Tous
les plis à $0^\circ$ : plaque très rigide selon $x$, molle selon $y$. Angles
bien répartis : plaque qui se comporte pareil dans toutes les directions.
Il y a $180^{20} \approx 1{,}3\times10^{45}$ empilements possibles :
impossible de tous les essayer, d'où l'optimisation.
\end{idee}

% =============================================================================
\section{La rigidité : du pli à la plaque}
% =============================================================================
\subsection{Loi de Hooke en contraintes planes}
Pour un ressort : force = raideur $\times$ allongement. Pour un matériau,
c'est pareil avec des \textbf{contraintes} $\sigma$ (force par surface, en
MPa) et des \textbf{déformations} $\varepsilon$ (allongement relatif, sans
unité). Une plaque est fine : on ne regarde que ce qui se passe \textbf{dans
son plan} (hypothèse des \emph{contraintes planes}). Il reste 3 contraintes
et 3 déformations : $\sigma_{11}, \sigma_{22}, \sigma_{12}$ et
$\varepsilon_{11}, \varepsilon_{22}, \gamma_{12} = 2\varepsilon_{12}$
(distorsion d'angle). On les relie par une matrice $3\times3$.

\subsection{La souplesse \texorpdfstring{$\Sm$}{S} et la rigidité \texorpdfstring{$\Q$}{Q}}
\emph{Fichier \code{laminate.py}, fonctions \code{ply\_compliance} et
\code{ply\_stiffness} (énoncé §2.b).}

Le plus naturel est d'écrire \og déformation = $\Sm$ $\times$ contrainte\fg,
car $\Sm$ se remplit directement avec les constantes de l'énoncé :
\[
\begin{Bmatrix}\varepsilon_{11}\\ \varepsilon_{22}\\ \gamma_{12}\end{Bmatrix}
=
\underbrace{\begin{bmatrix}
 \dfrac{1}{E_1} & -\dfrac{\nu_{12}}{E_1} & 0\\[1.2ex]
 -\dfrac{\nu_{12}}{E_1} & \dfrac{1}{E_2} & 0\\[1.2ex]
 0 & 0 & \dfrac{1}{G_{12}}
\end{bmatrix}}_{\Sm}
\begin{Bmatrix}\sigma_{11}\\ \sigma_{22}\\ \sigma_{12}\end{Bmatrix}
\]
Lecture : si je tire avec $\sigma_{11}$, le pli s'allonge de
$\sigma_{11}/E_1$ selon 1 et rétrécit de $\nu_{12}\sigma_{11}/E_1$ selon 2.
Les zéros disent que tirer ne crée pas de cisaillement (orthotropie, dans les
axes du pli). On veut plutôt \og contrainte = $\Q$ $\times$ déformation\fg :
\[
\boxed{\;\Q = \Sm^{-1}\;}
\qquad\text{(matrice de rigidité réduite du pli, l'objet central du projet)}
\]

\begin{idee}[Les indices 1, 2, 6 des matrices]
Les cases d'une matrice $3\times3$ comme $\Q$ sont numérotées avec la
\og notation de Voigt\fg{} : l'indice \textbf{1} correspond à la direction
11 (traction selon 1), \textbf{2} à la direction 22 (traction selon 2) et
\textbf{6} au cisaillement 12. Ainsi :
\[
\Q =
\begin{bmatrix}
Q_{11} & Q_{12} & Q_{16}\\
Q_{12} & Q_{22} & Q_{26}\\
Q_{16} & Q_{26} & Q_{66}
\end{bmatrix}
\qquad
\begin{array}{l}
Q_{11} : \text{rigidité en traction selon 1}\\
Q_{22} : \text{rigidité en traction selon 2}\\
Q_{66} : \text{rigidité en cisaillement}\\
Q_{12} : \text{effet Poisson (tirer selon 1 charge selon 2)}\\
Q_{16}, Q_{26} : \text{couplage traction / cisaillement}
\end{array}
\]
Dans les axes du pli, $Q_{16} = Q_{26} = 0$ (orthotropie). Ils deviennent
non nuls dès qu'on tourne le pli : tirer en biais sur des fibres les fait
\og glisser\fg.
\end{idee}

\subsection{Tourner un pli : \texorpdfstring{$\Q(\delta)$}{Q(delta)}}
\emph{Fonction \code{rotate\_stiffness} (P2 slide 27).}

$\Q$ est écrite dans les axes des fibres. Pour un pli tourné de $\delta$, il
faut l'exprimer dans les axes $x, y$ de la plaque. Avec $c = \cos\delta$ et
$s = \sin\delta$, par exemple :
\[
\bar Q_{11} = Q_{11}c^4 + 2(Q_{12}+2Q_{66})s^2c^2 + Q_{22}s^4,
\qquad
\bar Q_{16} = (Q_{11}-Q_{12}-2Q_{66})\,s c^3 + (Q_{12}-Q_{22}+2Q_{66})\,s^3 c,
\]
et quatre autres formules du même genre. C'est lourd : c'est justement pour
ça que la méthode polaire (section 4) existe.

\subsection{Membrane, flexion et couplage : \texorpdfstring{$\mathbf{A}, \mathbf{B}, \mathbf{D}$}{A, B, D}}
\emph{Fonction \code{abd\_matrices} (P2 slide 41).}

Une plaque peut être sollicitée de deux façons :
\begin{itemize}
  \item \textbf{en membrane} : on la tire ou la cisaille dans son plan (une
        feuille qu'on étire). Efforts $\mathbf{N}$, déformations du plan
        moyen $\boldsymbol{\varepsilon}_0$ ;
  \item \textbf{en flexion} : on la plie (une règle qu'on courbe). Moments
        $\mathbf{M}$, courbures $\boldsymbol{\chi}$.
\end{itemize}

\begin{center}
\begin{tikzpicture}[>={Stealth[length=2.4mm]}]
  % ---------------- (a) Membrane ----------------
  \begin{scope}
    % plate before (dashed) and after stretching
    \draw[dashed,gray] (0.3,0) rectangle (3.7,0.3);
    \fill[blue!12] (0,0) rectangle (4,0.3);
    \draw (0,0) rectangle (4,0.3);
    \draw[dashed,red!60!black] (0,0.15) -- (4,0.15);
    % forces N
    \draw[->,very thick] (0,0.15) -- (-0.9,0.15) node[left] {$\mathbf{N}$};
    \draw[->,very thick] (4,0.15) -- (4.9,0.15) node[right] {$\mathbf{N}$};
    \node[align=center,font=\small] at (2,-0.75)
      {(a) \textbf{Membrane} : on tire dans le plan,\\la plaque s'allonge et reste plate\\(rigidité $\mathbf{A}$, déformation $\boldsymbol{\varepsilon}_0$)};
  \end{scope}
  % ---------------- (b) Flexion ----------------
  \begin{scope}[xshift=9.6cm]
    % bent plate: band between two concentric arcs
    \fill[blue!12] (-1.9,0.55) arc[start angle=118,end angle=62,radius=4.0]
                  -- ++(0,0) -- (1.75,0.25) arc[start angle=62,end angle=118,radius=3.7] -- cycle;
    \draw (-1.9,0.55) arc[start angle=118,end angle=62,radius=4.0];
    \draw (-1.73,0.25) arc[start angle=118,end angle=62,radius=3.7];
    \draw (-1.9,0.55) -- (-1.73,0.25);
    \draw (1.9,0.55) -- (1.73,0.25);
    % moments M (curved arrows at both ends)
    \draw[->,very thick] (-2.15,-0.1) arc[start angle=220,end angle=110,radius=0.45];
    \node[left] at (-2.55,0.45) {$\mathbf{M}$};
    \draw[->,very thick] (2.15,-0.1) arc[start angle=-40,end angle=70,radius=0.45];
    \node[right] at (2.55,0.45) {$\mathbf{M}$};
    \node[align=center,font=\small] at (0,-0.75)
      {(b) \textbf{Flexion} : on plie la plaque,\\elle se courbe\\(rigidité $\mathbf{D}$, courbure $\boldsymbol{\chi}$)};
  \end{scope}
\end{tikzpicture}

\vspace{0.6em}
\begin{tikzpicture}[>={Stealth[length=2.4mm]}]
  % ---------------- (c) Coupling ----------------
  \fill[orange!15] (-1.9,0.55) arc[start angle=118,end angle=62,radius=4.0]
                -- (1.73,0.25) arc[start angle=62,end angle=118,radius=3.7] -- cycle;
  \draw (-1.9,0.55) arc[start angle=118,end angle=62,radius=4.0];
  \draw (-1.73,0.25) arc[start angle=118,end angle=62,radius=3.7];
  \draw (-1.9,0.55) -- (-1.73,0.25);
  \draw (1.9,0.55) -- (1.73,0.25);
  \draw[->,very thick] (-1.82,0.4) -- (-2.8,0.15) node[left] {$\mathbf{N}$};
  \draw[->,very thick] (1.82,0.4) -- (2.8,0.15) node[right] {$\mathbf{N}$};
  \node[align=center,font=\small] at (0,-0.6)
    {(c) \textbf{Couplage} ($\mathbf{B}\neq0$) : on ne fait que tirer\ldots{} et pourtant la plaque se courbe.\\
     C'est ce qu'on veut supprimer avec l'objectif $\Phi_1$ ($\Bs = 0$).};
\end{tikzpicture}
\end{center}
Les dessins sont des vues de côté (la plaque est vue sur la tranche,
l'épaisseur est très exagérée). La ligne pointillée rouge en (a) est le
\textbf{plan moyen}, à mi-épaisseur ($z = 0$).
La Théorie Classique des Stratifiés relie tout ça :
\[
\mathbf{N} = \mathbf{A}\,\boldsymbol{\varepsilon}_0 + \mathbf{B}\,\boldsymbol{\chi},
\qquad
\mathbf{M} = \mathbf{B}\,\boldsymbol{\varepsilon}_0 + \mathbf{D}\,\boldsymbol{\chi}.
\]
$\mathbf{A}$ : rigidité de \textbf{membrane} ; $\mathbf{D}$ : rigidité de
\textbf{flexion} ; $\mathbf{B}$ : \textbf{couplage}. Si $\mathbf{B}\neq0$,
tirer sur la plaque la fait aussi se courber (et inversement) : une pièce qui
se tord quand on la tire, ou qui gauchit en refroidissant après cuisson.

Le pli $k$ occupe $z_{k-1} \le z \le z_k$, avec $z=0$ au milieu,
$z_0 = -h/2$ en bas et $z_n = +h/2$ en haut :
\[
\mathbf{A} = \sum_{k=1}^{n} \Q_k\,(z_k - z_{k-1}),
\qquad
\mathbf{B} = \frac12\sum_{k=1}^{n} \Q_k\,(z_k^2 - z_{k-1}^2),
\qquad
\mathbf{D} = \frac13\sum_{k=1}^{n} \Q_k\,(z_k^3 - z_{k-1}^3).
\]

\begin{center}
\begin{tikzpicture}[scale=0.9]
  \foreach \i/\c in {0/blue!15,1/blue!25,2/blue!15,3/blue!25,4/blue!15,5/blue!25}{
    \fill[\c] (0,\i*0.5) rectangle (6,\i*0.5+0.5);
    \draw (0,\i*0.5) rectangle (6,\i*0.5+0.5);}
  \draw[dashed,red!70!black,thick] (-0.4,1.5) -- (6.4,1.5) node[right] {$z=0$ (plan moyen)};
  \node[left] at (0,0) {$z_0=-h/2$};
  \node[left] at (0,3) {$z_n=+h/2$};
  \node at (3,0.25) {\small pli 1};
  \node at (3,2.75) {\small pli $n$};
  \draw[-{Stealth}] (7.2,0) -- (7.2,3.2) node[above] {$z$};
\end{tikzpicture}
\end{center}

\begin{idee}[Intuition]
\begin{itemize}
  \item $\mathbf{A}$ est une simple \textbf{moyenne} des plis : l'ordre
        d'empilement ne compte pas.
  \item $\mathbf{D}$ pondère par $z^3$ : les plis \textbf{à l'extérieur}
        comptent beaucoup plus que ceux du milieu (comme une poutre en I).
  \item $\mathbf{B}$ pondère par $z^2$ avec un signe : un pli en bas et le
        même pli en haut s'annulent. Un stratifié \textbf{symétrique}
        (miroir par rapport au milieu, ex. $[0/45/45/0]$) a donc toujours
        $\mathbf{B}=0$.
\end{itemize}
\end{idee}

\subsection{Matrices homogénéisées}
\emph{Fonction \code{homogenised\_matrices}.}

$\mathbf{A}$, $\mathbf{B}$, $\mathbf{D}$ n'ont pas les mêmes unités. On les
normalise pour les comparer au pli (toutes deviennent des MPa) :
\[
\boxed{\;\As = \frac{\mathbf{A}}{h},\qquad
\Bs = \frac{2\,\mathbf{B}}{h^2},\qquad
\Ds = \frac{12\,\mathbf{D}}{h^3},\qquad
\Cm = \As - \Ds\;}
\]
Exemple : tous les plis à $0^\circ$ donnent $\As = \Ds = \Q$ et $\Bs = 0$.
$\Cm$ mesure la différence entre membrane et flexion ($\Cm=0$ : même
comportement dans les deux cas, on dit \emph{homogène}).

% =============================================================================
\section{La méthode polaire (Verchery, 1979)}
% =============================================================================
\emph{Fichier \code{polar.py} (P2 slides 35 à 42).}

\subsection{Le problème des composantes cartésiennes}
Une matrice comme $\Q$ a 6 nombres utiles, mais ils \textbf{changent quand
on tourne les axes}. Difficile de répondre à \og cette plaque est-elle
isotrope ?\fg{} en regardant juste les composantes.

\subsection{Les six paramètres polaires}

\begin{idee}[Qui est \texorpdfstring{$\mathbf{L}$}{L} ?]
Dans cette section, $\mathbf{L}$ n'est \textbf{pas une nouvelle matrice} :
c'est un \textbf{nom générique} pour \og n'importe laquelle de nos matrices
de rigidité\fg. La méthode polaire marche exactement de la même façon pour
toutes : on peut remplacer $\mathbf{L}$ par
\[
\Q\ \text{(le pli)},\qquad \As\ \text{(membrane)},\qquad \Bs\ \text{(couplage)},
\qquad \Ds\ \text{(flexion)},\qquad \Cm = \As - \Ds .
\]
Par exemple, avec $\mathbf{L} = \Q$, la formule de $T_0$ ci-dessous devient
$8\,T_0 = Q_{11} - 2Q_{12} + 4Q_{66} + Q_{22}$. Avec $\mathbf{L} = \As$,
elle donne le $T_0$ de la membrane, noté $T_0^{A^*}$, et ainsi de suite.
C'est d'ailleurs ce que fait le code : une seule fonction,
\code{polar\_parameters(L)}, est appelée cinq fois, une fois par matrice.

Le cours écrit les composantes avec quatre indices ($L_{1111}$, $L_{1122}$,
$L_{1212}$\ldots) parce qu'une rigidité est, en toute rigueur, un tenseur
d'ordre 4. Ce sont les mêmes cases que celles de la matrice $3\times3$, avec
la correspondance :
\[
L_{1111} = L_{11},\quad L_{2222} = L_{22},\quad L_{1122} = L_{12},\quad
L_{1212} = L_{66},\quad L_{1112} = L_{16},\quad L_{2212} = L_{26}.
\]
\end{idee}

On réécrit donc chacune de ces matrices avec six autres nombres, mieux
choisis :
\begin{center}
\begin{tabularx}{\linewidth}{l l X}
\toprule
Paramètres & Nom & Sens physique\\
\midrule
$T_0, T_1$ & modules isotropes & la partie identique dans toutes les directions\\
$R_0, R_1$ & modules anisotropes & \textbf{combien} la rigidité varie avec la direction\\
$\Phi_0, \Phi_1$ & angles polaires & \textbf{dans quelle direction} sont les maxima\\
\bottomrule
\end{tabularx}
\end{center}
Formules (slide 36), écrites avec les indices de la matrice $3\times3$ :
\begin{align*}
8\,T_0 &= L_{11} - 2L_{12} + 4L_{66} + L_{22}\\
8\,T_1 &= L_{11} + 2L_{12} + L_{22}\\
8\,R_0\,e^{4i\Phi_0} &= L_{11} + 4i\,L_{16} - 2L_{12} - 4L_{66} - 4i\,L_{26} + L_{22}\\
8\,R_1\,e^{2i\Phi_1} &= L_{11} + 2i\,L_{16} + 2i\,L_{26} - L_{22}
\end{align*}
On utilise des nombres complexes : $R$ est le module, $\Phi$ vient de
l'argument. Et dans l'autre sens (slide 35), par exemple :
\[
L_{11} = T_0 + 2T_1 + R_0\cos 4\Phi_0 + 4R_1\cos 2\Phi_1 .
\]

\begin{idee}[Pourquoi c'est génial]
\begin{enumerate}
  \item $T_0, T_1, R_0, R_1$ \textbf{ne changent pas quand on tourne les
        axes} (ce sont des \emph{invariants}). Tourner les axes de $\theta$
        revient juste à faire $\Phi_0 \to \Phi_0-\theta$ et
        $\Phi_1 \to \Phi_1-\theta$ (slide 37).
  \item \textbf{Les symétries se lisent directement} sur les paramètres
        (tableau ci-dessous).
\end{enumerate}
\end{idee}

\begin{center}
\begin{tabularx}{\linewidth}{l l X}
\toprule
Symétrie (slide 38) & Condition polaire & Forme du diagramme polaire\\
\midrule
Orthotropie ordinaire & $\Phi_0-\Phi_1 = K\,\pi/4$, $K\in\{0,1\}$ & \og os\fg{} ou \og nœud papillon\fg\\
\textbf{Orthotropie $R_0$} & $\mathbf{R_0 = 0}$ & \og cacahuète\fg{} (2 lobes)\\
\textbf{Symétrie du carré} & $\mathbf{R_1 = 0}$ & \og carré arrondi\fg{} (4 lobes)\\
Isotropie & $R_0 = R_1 = 0$ & cercle parfait\\
\bottomrule
\end{tabularx}
\end{center}
Le \textbf{diagramme polaire} trace la rigidité $L_{11}(\theta)$ dans chaque
direction $\theta$ sur un graphique en rond : c'est la rigidité en traction
quand on tire dans la direction $\theta$. C'est ce que fait
\code{plots.py}.

\subsection{Le stratifié en polaire : des sommes de flèches}
C'est \textbf{le} résultat à comprendre (slides 41--42). Pour des plis
identiques :
\begin{itemize}
  \item $T_0$ et $T_1$ de $\As$ et $\Ds$ \textbf{sont ceux du pli} : les
        angles ne les changent pas, ils ne servent à rien dans
        l'optimisation ;
  \item $\Bs$ et $\Cm$ ont $T_0 = T_1 = 0$ ;
  \item tout ce qu'on contrôle, ce sont les $R$ et les $\Phi$ :
\end{itemize}
\begin{align*}
\As:\quad R_0^{A^*} e^{4i\Phi_0^{A^*}} &= \frac{1}{n}\,R_0\,e^{4i\Phi_0}\sum_{k=1}^{n} e^{4i\delta_k},
&
R_1^{A^*} e^{2i\Phi_1^{A^*}} &= \frac{1}{n}\,R_1\,e^{2i\Phi_1}\sum_{k=1}^{n} e^{2i\delta_k},\\
\Bs:\quad R_0^{B^*} e^{4i\Phi_0^{B^*}} &= \frac{1}{n^2}\,R_0\,e^{4i\Phi_0}\sum_{k=1}^{n} b_k\,e^{4i\delta_k},
&
R_1^{B^*} e^{2i\Phi_1^{B^*}} &= \frac{1}{n^2}\,R_1\,e^{2i\Phi_1}\sum_{k=1}^{n} b_k\,e^{2i\delta_k},\\
\Ds:\quad R_0^{D^*} e^{4i\Phi_0^{D^*}} &= \frac{1}{n^3}\,R_0\,e^{4i\Phi_0}\sum_{k=1}^{n} d_k\,e^{4i\delta_k},
&
R_1^{D^*} e^{2i\Phi_1^{D^*}} &= \frac{1}{n^3}\,R_1\,e^{2i\Phi_1}\sum_{k=1}^{n} d_k\,e^{2i\delta_k},
\end{align*}
avec les poids de position (déduits de $\Bs = 2\mathbf{B}/h^2$ et
$\Ds = 12\mathbf{D}/h^3$) :
\[
b_k = 2k - n - 1,
\qquad
d_k = 12\,k\,(k-n-1) + 4 + 3n(n+2).
\]
Les $R$, $\Phi$ sans exposant sont ceux du pli (ici $\Phi_0=\Phi_1=0$).

\begin{idee}[Image mentale : additionner des flèches]
Chaque pli est une \textbf{flèche} de longueur 1 qui pointe dans la
direction $4\delta_k$ (pour $R_0$) ou $2\delta_k$ (pour $R_1$). On additionne
les flèches bout à bout. \textbf{Si elles se referment sur elles-mêmes, la
somme est nulle et $R = 0$.}
\end{idee}

\begin{center}
\begin{tikzpicture}[>={Stealth[length=2.5mm]},thick]
  % [0/90] for R1: angles 2*delta = 0 and 180
  \begin{scope}
    \draw[->,blue!70!black] (0,0) -- (1.6,0);
    \draw[->,red!70!black] (1.6,0.08) -- (0,0.08);
    \node[below,align=center] at (0.8,-0.5) {\small $[0/90]$, flèches $2\delta$ :\\ \small $0^\circ$ et $180^\circ$ $\Rightarrow R_1^{A^*}=0$};
  \end{scope}
  % [0/45] for R0: angles 4*delta = 0 and 180
  \begin{scope}[xshift=5cm]
    \draw[->,blue!70!black] (0,0) -- (1.6,0);
    \draw[->,red!70!black] (1.6,0.08) -- (0,0.08);
    \node[below,align=center] at (0.8,-0.5) {\small $[0/45]$, flèches $4\delta$ :\\ \small $0^\circ$ et $180^\circ$ $\Rightarrow R_0^{A^*}=0$};
  \end{scope}
  % [0/60/-60]: 3 arrows at 120 deg forming a triangle
  \begin{scope}[xshift=10cm]
    \draw[->,blue!70!black] (0,0) -- (1.6,0);
    \draw[->,red!70!black] (1.6,0) -- ++(120:1.6);
    \draw[->,green!50!black] ($(1.6,0)+(120:1.6)$) -- (0,0);
    \node[below,align=center] at (0.8,-0.5) {\small $[0/60/{-60}]$ : triangle fermé\\ \small $\Rightarrow R_0^{A^*}=R_1^{A^*}=0$ (isotrope)};
  \end{scope}
\end{tikzpicture}
\end{center}

Pour $\Bs$ et $\Ds$, chaque flèche est en plus multipliée par son poids
$b_k$ ou $d_k$ (négatif en bas pour $b_k$, grand à l'extérieur pour $d_k$).
C'est pour ça qu'il est difficile d'annuler plusieurs choses à la fois : les
\textbf{mêmes} 20 angles doivent refermer plusieurs sommes de flèches avec
des poids différents.

\subsection{La norme d'un tenseur}
Pour mesurer la \og taille\fg{} d'un tenseur (slide 45) :
\[
\|\mathbf{L}\| = \sqrt{T_0^2 + 2\,T_1^2 + R_0^2 + 4\,R_1^2}.
\]
Pour notre pli, $\|\Q\| \approx 57\,400$ MPa.

% =============================================================================
\section{Ce que doit faire notre plaque (groupe 1)}
% =============================================================================
\emph{Fichier \code{fitness.py} (P2 slides 44 à 47), données dans
\code{config.py}.}

L'énoncé demande trois symétries. Chacune devient un \textbf{objectif
partiel} $\Phi_i \ge 0$ qui vaut \textbf{0 quand la symétrie est
parfaitement atteinte} :
\begin{center}
\begin{tabularx}{\linewidth}{X l l c}
\toprule
Symétrie demandée & Condition & Objectif & Formule\\
\midrule
Découplage & $\Bs = 0$ & $\Phi_1$ & $\left(\dfrac{\|\Bs\|}{\|\Q\|}\right)^{2}$\\[2.2ex]
Orthotropie $R_0$ en membrane & $R_0^{A^*}=0$ & $\Phi_4$ & $\left(\dfrac{R_0^{A^*}}{R_0}\right)^{2}$\\[2.2ex]
Symétrie du carré en flexion & $R_1^{D^*}=0$ & $\Phi_9$ & $\left(\dfrac{R_1^{D^*}}{R_1}\right)^{2}$\\[1.5ex]
\bottomrule
\end{tabularx}
\end{center}
On divise par une grandeur du pli pour que chaque $\Phi_i$ soit \textbf{sans
unité} et de l'ordre de 1 au pire : on peut alors les additionner. La
\textbf{fonction de coût} est la somme (slide 44) :
\[
\boxed{\;\Phi(\boldsymbol\delta) = \Phi_1(\boldsymbol\delta) + \Phi_4(\boldsymbol\delta) + \Phi_9(\boldsymbol\delta),
\qquad \min \Phi = 0\;}
\qquad \boldsymbol\delta = (\delta_1, \dots, \delta_{20}).
\]
Toute la difficulté : $\Phi$ est très \textbf{non convexe} en fonction des
angles (des milliers de creux locaux). Les méthodes classiques (descente de
gradient) y restent coincées, d'où l'algorithme génétique.

\begin{attention}[Deux points tranchés par le groupe]
\textbf{1. $\Phi_1$ et la slide 51.} Avec la formule du cours, l'exemple de
la slide 51 ($[9/0/18/{-4}/14/5]$, 6 plis T300/5208) donne
$\Phi = 9{,}96\times10^{-8}$. La slide affiche $3{,}5873\times10^{-6}$,
exactement $n^2 = 36$ fois plus : la slide utilise visiblement $n\,\Bs$ au
lieu de $\Bs$. On garde la formule du cours et on considère la valeur de la
slide comme une erreur d'application numérique.

\textbf{2. \og Orthotropie $R_1$\fg} n'existe pas dans le cours : c'est
une coquille de l'énoncé, confirmée par M. Montemurro. Il faut lire
\og symétrie du carré $R_1=0$ en flexion\fg, c'est-à-dire $R_1^{D^*}=0$
($\Phi_9$), ce que fait déjà le code.
\end{attention}

\begin{attention}[Pièges évités (à citer dans le rapport)]
\begin{itemize}
  \item Quand $R_0 = 0$ ou $R_1 = 0$, l'angle $\Phi_0$ ou $\Phi_1$
        \textbf{n'a plus de sens} (une flèche de longueur nulle n'a pas de
        direction). On ne combine donc jamais un objectif sur les angles
        ($\Phi_3$, $\Phi_7$) avec un objectif qui annule un $R$ sur le même
        tenseur ; \code{check\_objectives} refuse ces combinaisons.
  \item $\Phi_0$ est défini à $\pi/2$ près et $\Phi_1$ à $\pi$ près : pour
        $\Phi_3$ et $\Phi_7$, l'écart $\Phi_0-\Phi_1$ est ramené dans le bon
        intervalle avant d'être comparé à $K\pi/4$ (\code{wrap\_angle}).
\end{itemize}
Notre groupe n'utilise aucun objectif d'angle : pas de problème pour nous.
\end{attention}

% =============================================================================
\section{L'algorithme génétique (AG)}
% =============================================================================
\emph{Fichier \code{ga.py} (P1 slide 48, P2 slide 50).}

\subsection{L'idée}
On imite Darwin : une \textbf{population} de solutions ; les meilleures ont
plus d'\textbf{enfants} ; les enfants \textbf{mélangent} les caractères de
leurs parents ; de petites \textbf{mutations} apportent de la nouveauté. Au
fil des générations, la population s'améliore.
\begin{center}
\begin{tabular}{l l}
\toprule
Biologie & Notre problème\\
\midrule
individu & un stratifié (20 angles)\\
gène & un angle de pli\\
génotype (ADN) & les 160 bits qui codent les 20 angles\\
phénotype & les 20 angles en degrés\\
adaptation au milieu & fitness, calculée à partir du coût $\Phi$\\
\bottomrule
\end{tabular}
\end{center}

\subsection{Le codage binaire et les 76 codes en trop}
Chaque angle est stocké sur \textbf{8 bits} : $2^7 = 128$ valeurs ne
suffisent pas pour 180 angles, $2^8 = 256$ suffisent. Un individu fait donc
$20\times8 = 160$ bits. Mais $256 - 180 = 76$ codes ne correspondent à aucun
angle. Deux solutions :
\begin{itemize}
  \item \textbf{modulo} ($\text{code} \bmod 180$) : simple, mais les angles
        $-89^\circ$ à $-14^\circ$ auraient deux codes chacun, donc
        \textbf{deux fois plus de chances} d'apparaître : recherche biaisée ;
  \item \textbf{rejet} (notre choix) :
        \[
        \delta = -89^\circ + \text{code},\qquad \text{code}\in\{0,\dots,179\},
        \]
        et les codes 180 à 255 sont interdits. La population de départ n'a
        que des codes valides ; si un croisement ou une mutation crée un code
        interdit, ce gène est \textbf{rejeté} et l'enfant garde le gène de
        son parent. Chaque angle a la même chance ($1/180$). Sur un calcul
        complet, environ 2\,500 gènes sont rejetés en 500 générations
        (colonne \code{rejected\_genes} du CSV).
\end{itemize}

\subsection{Les étapes}
Les étapes sont décrites une par une, puis résumées par le schéma de la
boucle à la fin de la section.

\paragraph{1. Population initiale.} $N = 500$ stratifiés tirés au hasard.

\paragraph{2. Évaluation.} Pour chaque individu : décoder les angles,
calculer $\As, \Bs, \Ds$, leurs paramètres polaires, puis
$\Phi = \Phi_1 + \Phi_4 + \Phi_9$.

\paragraph{3. Fitness} (énoncé §2.a, avec $p = 1$) :
\[
f_i = \left(1 - \frac{\Phi_i - \min_j \Phi_j}{\max_j \Phi_j - \min_j \Phi_j}\right)^{p}.
\]
Le meilleur de la génération a $f = 1$, le pire $f = 0$. Exemple avec trois
individus de coûts $0{,}01$ ; $0{,}05$ ; $0{,}11$ : fitness $1$ ; $0{,}6$ ;
$0$.

\paragraph{4. Sélection par roulette biaisée.} Chaque individu reçoit une
part de roue proportionnelle à sa fitness :
$P_i = f_i / \sum_j f_j$. On tourne la roue $N$ fois pour choisir $N$
parents (un bon individu peut sortir plusieurs fois). Dans l'exemple :
$1/1{,}6 = 62{,}5\,\%$, $0{,}6/1{,}6 = 37{,}5\,\%$ et $0\,\%$.

\paragraph{5. Croisement} (probabilité $p_c = 0{,}85$). Les parents vont par
deux ; on coupe les deux chaînes de bits au même endroit et on échange les
fins :
\[
\begin{array}{lcl}
\text{parent 1 : } \texttt{10110|010} & \longrightarrow & \text{enfant 1 : } \texttt{10110|111}\\
\text{parent 2 : } \texttt{01100|111} & \longrightarrow & \text{enfant 2 : } \texttt{01100|010}
\end{array}
\]
Dans 15\,\% des cas, pas de croisement : les enfants sont des copies.

\paragraph{6. Mutation} (probabilité $p_m = 1/N$ par gène, comme dans le
cours). Chaque gène a une petite chance de voir un de ses 8 bits inversé
($0 \leftrightarrow 1$). Ça évite que toute la population devienne identique
et reste coincée.

\paragraph{7. Élitisme.} Le meilleur individu de la génération précédente
remplace le pire enfant : la meilleure solution trouvée \textbf{n'est jamais
perdue}. C'est pour ça que la courbe du meilleur coût ne remonte jamais.

\paragraph{8. Arrêt} après $N_{\text{gen}} = 500$ générations.

\paragraph{Résumé de la boucle.}
\begin{center}
\begin{tikzpicture}[
  node distance=0.55cm,
  box/.style={rectangle,draw,rounded corners,fill=blue!6,minimum width=6.4cm,minimum height=0.75cm,align=center,font=\small},
  dec/.style={diamond,draw,fill=orange!10,aspect=2.6,inner sep=1pt,font=\small,align=center},
  >={Stealth}]
  \node[box,fill=green!8] (init) {Population initiale aléatoire ($N$ individus)};
  \node[box,below=of init] (eval) {Évaluation : coût $\Phi$ de chaque individu};
  \node[box,below=of eval] (fit) {Fitness $f_i \in [0,1]$};
  \node[box,below=of fit] (sel) {Sélection par roulette $\rightarrow$ parents};
  \node[box,below=of sel] (cross) {Croisement (probabilité $p_c$) $\rightarrow$ enfants};
  \node[box,below=of cross] (mut) {Mutation (probabilité $p_m$ par gène)};
  \node[box,below=of mut] (rej) {Rejet des codes invalides};
  \node[box,below=of rej] (eli) {Élitisme : le meilleur parent remplace le pire enfant};
  \node[dec,below=0.7cm of eli] (stop) {génération\\$= N_{\text{gen}}$ ?};
  \node[box,fill=green!8,below=0.7cm of stop] (end) {Meilleur individu $\rightarrow$ résultats};
  \foreach \a/\b in {init/eval,eval/fit,fit/sel,sel/cross,cross/mut,mut/rej,rej/eli,eli/stop}
    \draw[->] (\a) -- (\b);
  \draw[->] (stop) -- node[right,font=\small]{oui} (end);
  \draw[->] (stop.east) -- node[above,font=\small]{non} ++(2.6,0) |- (eval.east);
\end{tikzpicture}
\end{center}

\subsection{Les paramètres (\code{config.py})}
\begin{center}
\begin{tabular}{l l l}
\toprule
Paramètre & Valeur & Rôle\\
\midrule
$N$ & 500 & taille de la population\\
$N_{\text{gen}}$ & 500 & nombre de générations\\
$p_c$ & 0,85 & probabilité de croisement\\
$p_m$ & $1/N = 0{,}002$ & probabilité de mutation d'un gène\\
$p$ & 1 & pression de sélection (exposant de la fitness)\\
seed & 1 & graine du hasard : même graine $\Rightarrow$ même résultat\\
\bottomrule
\end{tabular}
\end{center}
Ce sont les valeurs de la slide 50 du cours.

% =============================================================================
\section{Le code, fichier par fichier}
% =============================================================================
\begin{center}
\begin{tabularx}{\linewidth}{l X l}
\toprule
Fichier & Contenu & Cours\\
\midrule
\code{config.py} & toutes les données : pli, groupe, objectifs, paramètres AG & énoncé §1, P2 sl.~50\\
\code{laminate.py} & $\Sm$, $\Q$, $\Q(\delta)$, $\mathbf{A},\mathbf{B},\mathbf{D}$, $\As,\Bs,\Ds,\Cm$ & P2 sl.~27, 41\\
\code{polar.py} & $T_0,T_1,R_0,R_1,\Phi_0,\Phi_1$, norme, sommes polaires directes & P2 sl.~35--42, 45\\
\code{fitness.py} & les 11 objectifs partiels, le coût $\Phi$, la fitness & P2 sl.~44--47\\
\code{ga.py} & codage, population, roulette, croisement, mutation, rejet, élitisme & P1 sl.~48, P2 sl.~50\\
\code{plots.py} & courbe de convergence, diagrammes polaires & énoncé §3\\
\code{main.py} & enchaîne tout et écrit \code{results/} & ---\\
\code{tests/} & vérifications automatiques & P2 sl.~49, 51\\
\bottomrule
\end{tabularx}
\end{center}

\paragraph{Chaîne d'appel lors d'un calcul.}
\begin{center}
\begin{tikzpicture}[
  f/.style={rectangle,draw,fill=gray!6,font=\small\ttfamily,anchor=west,inner sep=3pt},
  n/.style={font=\small,anchor=west,text=gray!40!black},
  >={Stealth}]
  % Each call is a box; indentation = depth of the call.
  % Arrows leave from the left part of the caller and enter the callee.
  \node[f] (main) at (0,0)     {main.py};
  \node[f] (ply)  at (1,-0.8)  {laminate.ply\_stiffness};
  \node[f] (ga)   at (1,-1.6)  {ga.run\_ga};
  \node[f] (cost) at (2,-2.4)  {fitness.cost};
  \node[f] (lam)  at (3,-3.2)  {laminate.laminate\_matrices};
  \node[f] (pol)  at (3,-4.0)  {polar.polar\_parameters};
  \node[f] (obj)  at (3,-4.8)  {fitness.partial\_objectives};
  \node[f] (out)  at (1,-5.6)  {main: CSV, best\_laminate.txt, plots};

  \node[n] at (8.6,-0.8) {$\Sm$, $\Q$ du pli};
  \node[n] at (8.6,-1.6) {boucle sur les générations};
  \node[n] at (8.6,-2.4) {pour chaque individu};
  \node[n] at (8.6,-3.2) {$\As,\Bs,\Ds,\Cm$};
  \node[n] at (8.6,-4.0) {$T_0,T_1,R_0,R_1,\Phi_0,\Phi_1$};
  \node[n] at (8.6,-4.8) {$\Phi_1,\dots,\Phi_{11}$ (on garde 1, 4, 9)};
  \node[n] at (8.6,-5.6) {à la fin du calcul};

  \draw[->] ([xshift=4mm]main.south west) |- (ply.west);
  \draw[->] ([xshift=4mm]main.south west) |- (ga.west);
  \draw[->] ([xshift=4mm]main.south west) |- (out.west);
  \draw[->] ([xshift=4mm]ga.south west)   |- (cost.west);
  \draw[->] ([xshift=4mm]cost.south west) |- (lam.west);
  \draw[->] ([xshift=4mm]cost.south west) |- (pol.west);
  \draw[->] ([xshift=4mm]cost.south west) |- (obj.west);
\end{tikzpicture}
\end{center}

\paragraph{Comment on sait que la mécanique est juste.} Les tests
(\code{tests/test\_mechanics.py}) comparent le code aux valeurs du cours :
\begin{itemize}
  \item pli T300/5208 (slide 49) : $T_0 = 26\,880$, $T_1 = 24\,744$,
        $R_0 = 19\,710$, $R_1 = 21\,433$ MPa : \textbf{retrouvés} ;
  \item stratifié $[9/0/18/{-4}/14/5]$ (slide 51) : tous les paramètres de
        $\As, \Bs, \Ds$ \textbf{retrouvés à 4 décimales}, avec
        $\Phi_0 = \Phi_1 = 7^\circ$ ;
  \item contrôle croisé : la voie cartésienne (matrices) et la voie polaire
        (sommes de flèches) donnent les mêmes nombres sur des empilements
        tirés au hasard.
\end{itemize}
Pour lancer les tests puis le calcul (depuis le dossier
\code{Claude Momon's}) :
\begin{center}
\code{../.venv/Scripts/python -m pytest}\qquad
\code{../.venv/Scripts/python main.py}
\end{center}

% =============================================================================
\section{Lire les résultats}
% =============================================================================
Un calcul complet prend environ 35 s. Résultat avec la graine 1 :
\[
[22/{-41}/80/{-15}/{-89}/42/{-61}/61/{-74}/83/{-72}/{-38}/67/{-52}/43/84/5/38/{-25}/{-82}]
\]
\[
\Phi = 4{,}1\times10^{-5}
\qquad(\Phi_1 = 2{,}5\times10^{-5},\ \Phi_4 = 6{,}3\times10^{-6},\ \Phi_9 = 9{,}5\times10^{-6}).
\]
Concrètement : $R_0^{A^*} = 44$ MPa au lieu de 17\,700 pour le pli seul
(divisé par 400) et $R_1^{D^*} = 59$ MPa au lieu de 19\,000. C'est un très
bon \textbf{pseudo-optimum} : pas exactement 0 (impossible avec des angles
entiers), mais très proche.

\begin{itemize}
  \item \code{convergence.*} : en haut, le meilleur coût (en escalier, il ne
        remonte jamais grâce à l'élitisme), la moyenne et le pire, en échelle
        logarithmique ; en bas, l'écart-type : grand au début (population
        variée), petit à la fin (population regroupée).
  \item \code{polar\_diagrams.*} : la rigidité dans chaque direction. Le pli
        $\Q$ fait un \og 8\fg{} allongé à $0^\circ$ ; $\As$ a la forme de
        cacahuète de l'orthotropie $R_0$ ; $\Ds$ ressemble à un carré arrondi
        ; $\Bs$ est près de 100 fois plus petit que $\As$ (il a sa propre
        échelle, en MPa).
  \item \code{ga\_statistics.csv} : une ligne par génération (min, moyenne,
        médiane, max, écart-type, gènes rejetés, meilleur empilement).
        S'ouvre directement dans Excel.
  \item \code{best\_laminate.txt} : l'empilement, les $\Phi_i$, tous les
        paramètres polaires et les matrices.
\end{itemize}

% =============================================================================
\section{Ce qu'il reste à faire pour le rendu}
% =============================================================================
\begin{enumerate}
  \item \textbf{Plusieurs calculs} avec des graines différentes
        (\code{--seed 2}, \code{3}\ldots) et éventuellement le croisement en
        deux points (\code{--crossover two\_point}), pour montrer que le
        résultat est robuste.
  \item \textbf{Le rapport} (énoncé §3) : théorie (sections 2 à 5 de ce
        document), méthodologie (sections 6 et 7, choix du rejet, décision
        sur $\Phi_1$), résultats (section 8 et figures).
  \item \textbf{Le diagramme de flux} du code : les deux schémas des
        sections 6.3 et 7 en sont une bonne base.
\end{enumerate}

% =============================================================================
\section*{Glossaire}
\addcontentsline{toc}{section}{Glossaire}
% =============================================================================
\begin{description}[style=nextline,leftmargin=1.2em,itemsep=0.1em]
  \item[Anisotrope] propriétés différentes selon la direction.
  \item[Orthotrope] anisotrope avec deux directions de symétrie perpendiculaires.
  \item[Isotrope] identique dans toutes les directions (comme l'acier).
  \item[Contraintes planes] on néglige ce qui se passe dans l'épaisseur.
  \item[Invariant] grandeur qui ne change pas quand on tourne les axes.
  \item[Couplage membrane/flexion] tirer fait courber ($\mathbf{B}\neq0$).
  \item[Quasi-homogène] $\Bs = 0$ et $\Cm = 0$ (même comportement partout).
  \item[Fonction de coût] nombre à minimiser ; ici 0 = objectif atteint.
  \item[Fitness] note d'adaptation entre 0 et 1 ; plus grand = meilleur.
  \item[Élitisme] on garde toujours le meilleur individu.
  \item[Pseudo-optimum] très bonne solution, pas forcément la meilleure absolue.
\end{description}

\end{document}
